% AI4SE Workshop @ KI 2026 — Position Paper
% Springer LNCS format, max 4 pages
\documentclass[runningheads]{llncs}

\usepackage[T1]{fontenc}
\usepackage{graphicx}
\usepackage{hyperref}
\usepackage{microtype}
\usepackage{booktabs}

\begin{document}

\title{From Stateless to Cognitive: Temporal Memory as the Missing Layer in Agentic AI}

\titlerunning{Temporal Memory in Agentic AI}

\author{Michael Banf, Johannes Kuhn, Slawomir Garcarz \and Max Becker}

\authorrunning{Banf et al.}

\institute{Perelyn GmbH, Munich, Germany\\
\email{michael.banf@perelyn.com}}

\maketitle

\begin{abstract}
LLM-based agents operate with no durable memory by default. Recent
architectures---extraction-based (Mem0), self-managing (Letta), and
graph-based (Graphiti)---address persistence, yet empirical evaluations
reveal a shared blind spot: \emph{temporal cognition}. Agents store facts
but cannot reliably distinguish current from outdated knowledge, track
belief evolution, or consolidate experiences across temporal distance. The
LoCoMo-Plus experiment finds cognitive accuracy below 25\% across all three
paradigms, with a 35--45 percentage-point drop from factual to cognitive
tasks. We survey the rapidly growing architectural landscape (from $\sim$25
papers in 2023 to 200+ in 2025, with 110+ submissions to the ICLR\,2026
MemAgents workshop), analyze benchmark evidence for the temporal deficit,
and identify four open challenges. We argue that closing the temporal gap
is the prerequisite for genuine cognitive continuity.

\keywords{Agentic Memory \and Temporal Reasoning \and Cognitive Continuity
\and Knowledge Graphs \and LLM Agents}
\end{abstract}

%%% ------------------------------------------------------------------
\section{Introduction}
%%% ------------------------------------------------------------------

The dominant LLM deployment pattern is stateless. Agentic systems break
this pattern but expose a deeper problem: without persistent memory,
agents cannot learn from experience or reason about change over time.
The field has responded rapidly---from $\sim$25 papers in 2023 to 200+ in
2025, with the ICLR\,2026 MemAgents workshop accepting 70 papers from 110+
submissions~\cite{memagents2026}---yet our classification of all 70
papers finds that 66\% propose new systems while only 4\% implement
temporal mechanisms beyond basic timestamps. Every major AI provider ships
memory (OpenAI since Feb\,2024, Anthropic Sep\,2025, Google Mar\,2026);
the open-source ecosystem has matured (Mem0: 47K stars; Graphiti: 45K;
Letta: 13K)~\cite{hu2025}.

The CoALA taxonomy~\cite{sumers2024} maps human memory types onto agent
components: \emph{working} (context window), \emph{semantic} (facts),
\emph{episodic} (experiences), \emph{procedural} (skills). Three
production paradigms address persistence: extraction-based
(Mem0~\cite{mem0_2025}), self-managing (MemGPT/Letta~\cite{packer2023}),
and graph-based temporal
(Zep/Graphiti~\cite{rasmussen2025})~\cite{hu2025,zhang2024survey}.

Yet benchmark evidence is sobering. LoCoMo-Plus~\cite{gurawaliya2026}
found peak cognitive accuracy of just 24.2\% (Letta), 22.2\% (Zep), and
18.0\% (Mem0)---while the same systems achieve 59--63\% on factual
recall, a \textbf{35--45~pp factual-to-cognitive gap}.
MemoryAgentBench~\cite{memoryagentbench2026} confirms: on multi-hop
conflict resolution, \emph{all} methods score below 7\%. This is not a
retrieval failure; it is a \emph{temporal cognition} failure.

%%% ------------------------------------------------------------------
\section{Temporal Memory Architectures}
%%% ------------------------------------------------------------------

\paragraph{Bi-temporal knowledge graphs.}
Graphiti~\cite{rasmussen2025} distinguishes \emph{event time} (when
something happened) from \emph{ingestion time} (when the system learned
it), with every edge carrying explicit validity intervals. On the Deep
Memory Retrieval benchmark, Zep reaches 94.8\% (vs.\ MemGPT 93.4\%);
on LongMemEval, it achieves 71.2\% with gpt-4o using only 1.6K context
tokens versus 115K for full-context baselines~\cite{rasmussen2025}.

\paragraph{Semantic timelines and hierarchical consolidation.}
TSM~\cite{tsm2026} identifies \emph{temporal inaccuracy} (memories indexed
by dialogue time, not event time) and \emph{temporal fragmentation} (loss
of durative state). It consolidates point-wise observations into durative
memory, achieving 74.8\% on LongMemEval\_S---with temporal-category gains
of +22.6 points over the A-MEM baseline. TiMem~\cite{timem2026}
organizes conversations through a Temporal Memory Tree, reaching 75.3\%
on LoCoMo (up from MemoryBank's 39.8\%) while reducing recalled memory
tokens by 52\%. GAM~\cite{gam2026} decouples episodic from semantic
consolidation via an Event Progression Graph, improving Temporal F1
by 18\% over Mem0.

\paragraph{Cognitive-science-inspired decay.}
MemoryBank~\cite{zhong2024} implements Ebbinghaus-curve forgetting; the
ACT-R-inspired architecture of~\cite{actr2025} integrates
activation-based recall directly into generation, making forgetting
intrinsic rather than post-hoc.

\paragraph{Emerging: RL-based memory management.}
A new direction replaces heuristic policies with learned ones:
Memory-R1~\cite{memoryr1_2025} trains a Memory Manager via RL for
ADD/UPDATE/DELETE operations; AgeMem~\cite{agemem2026} unifies LTM/STM
management through RL-optimized tool actions---but both remain
early-stage.

No single paradigm masters all competencies; the evidence points toward
hybrids combining temporal grounding, hierarchical consolidation,
principled decay, and learned policies.

%%% ------------------------------------------------------------------
\section{Benchmark Evidence for the Temporal Deficit}
%%% ------------------------------------------------------------------

LoCoMo~\cite{maharana2024}, LongMemEval~\cite{wu2025},
MemoryAgentBench~\cite{memoryagentbench2026}, and
MemoryArena~\cite{memoryarena2026} evaluate temporal competence---yet
LoCoMo is used by only 19\% of workshop papers; 43\% use custom
benchmarks, fragmenting evaluation.

LoCoMo-Plus~\cite{gurawaliya2026} provides the most fine-grained
analysis. At short distance (1--2 weeks), Letta and Zep tie at 27.5\%;
at long distance (3+ months), all converge to $\sim$20\%. Of 401
dialogues, 59.1\% were wrong across \emph{all three} architectures;
only 6.7\% correct by all---the deficit is structural, not
architecture-specific.

%%% ------------------------------------------------------------------
\section{Four Open Challenges}
%%% ------------------------------------------------------------------

\paragraph{C1: Temporal grounding.}
Of 70 workshop papers, 54\% have no temporal mechanism; of the 46\%
that do, the vast majority use only timestamps or ordering---zero
implement bi-temporal tracking or temporal consistency checking.
Graphiti and TSM offer foundations, but neither handles temporal
boundary uncertainty, conflicting evidence, or gradual state
transitions. A unified temporal ontology is missing.

\paragraph{C2: Cognitive consolidation.}
Human memory consolidates experiences into abstract representations.
Agent analogs exist (TiMem's hierarchical tree, Generative Agents'
reflection~\cite{park2023}) but operate on ad-hoc schedules. Only
31\% of workshop papers address consolidation; just 19\% implement
explicit forgetting. Procedural memory---how to perform tasks---appears
in only 24\% of papers despite being central to agent autonomy.

\paragraph{C3: Collaborative memory, governance, and security.}
The workshop reveals a single-agent bias: only 7\% address
collaborative/shared memory; 6\% track provenance; 4\% address
security; 3\% handle multi-modal memory; 1\% touch interoperability;
zero papers target GDPR-compliant deletion as a primary contribution.
Yet MINJA~\cite{minja2026} achieves 98\% injection success on shared
memory agents. MemArchitect~\cite{memarchitect2026} and
SSGM~\cite{ssgm2026} propose governance middleware;
\emph{mnemonic sovereignty}~\cite{mnemonic2026} addresses EU AI Act
compliance---but multi-user access control, memory portability, and
causal provenance (3\%) remain open.

%%% ------------------------------------------------------------------
\section{Research Agenda and Conclusion}
%%% ------------------------------------------------------------------

Four research questions structure our agenda:
\textbf{RQ1.}~What temporal ontology unifies event time, ingestion time,
validity intervals, and gradual transitions?
\textbf{RQ2.}~What consolidation mechanisms preserve temporal provenance
while reducing memory volume?
\textbf{RQ3.}~How can governance policies (including GDPR-compliant
deletion) be enforced across heterogeneous memory backends?
\textbf{RQ4.}~How should temporal benchmarks evolve beyond QA to evaluate
belief revision, proactive recall, and temporal planning?

Our workshop analysis reveals an \emph{episodic-temporal paradox}:
59\% of papers address episodic memory, yet only 4\% implement
temporal mechanisms beyond timestamps---episodes are treated as
unordered collections rather than temporally situated experiences.
The 35--45~pp factual-to-cognitive gap persists across all
architectures. Closing it is the prerequisite for cognitive
continuity---agents that maintain coherent, evolving understanding
across unbounded interaction horizons.

%%% ------------------------------------------------------------------
\begin{thebibliography}{20}
%%% ------------------------------------------------------------------

\bibitem{sumers2024}
Sumers, T.R., Yao, S., Narasimhan, K., Griffiths, T.L.: Cognitive
Architectures for Language Agents. In: TMLR (2024)

\bibitem{park2023}
Park, J.S., et al.: Generative Agents: Interactive Simulacra of Human
Behavior. In: UIST 2023

\bibitem{mem0_2025}
Chhikara, P., et al.: Mem0: Building Production-Ready AI Agents with
Scalable Long-Term Memory. In: ECAI 2025

\bibitem{packer2023}
Packer, C., et al.: MemGPT: Towards LLMs as Operating Systems.
arXiv:2310.08560 (2023)

\bibitem{rasmussen2025}
Rasmussen, P., et al.: Zep: A Temporal Knowledge Graph Architecture for
Agent Memory. arXiv:2501.13956 (2025)

\bibitem{hu2025}
Hu, Y., et al.: Memory in the Age of AI Agents: A Survey---Forms,
Functions and Dynamics. arXiv:2512.13564 (2025)

\bibitem{zhang2024survey}
Zhang, Z., et al.: A Survey on the Memory Mechanism of Large Language
Model based Agents. In: ACM TOIS (2024)

\bibitem{gurawaliya2026}
Gurawaliya, H., Nalenz, M., Ma, M.: Agentic Memory Realized: From
Stateless Models to Cognitive Continuity. appliedAI Initiative (2026)

\bibitem{tsm2026}
Beyond Dialogue Time: Temporal Semantic Memory for Personalized LLM Agents.
arXiv:2601.07468 (2026)

\bibitem{timem2026}
TiMem: Temporal-Hierarchical Memory Consolidation for Long-Horizon
Conversational Agents. arXiv:2601.02845 (2026)

\bibitem{zhong2024}
Zhong, W., et al.: MemoryBank: Enhancing Large Language Models with
Long-Term Memory. In: AAAI 2024

\bibitem{actr2025}
Human-Like Remembering and Forgetting in LLM Agents: An ACT-R-Inspired
Memory Architecture. In: HAI~'25

\bibitem{maharana2024}
Maharana, A., et al.: Evaluating Very Long-Term Conversational Memory of
LLM Agents. In: ACL 2024

\bibitem{wu2025}
Wu, D., et al.: LongMemEval: Benchmarking Chat Assistants on Long-Term
Interactive Memory. In: ICLR 2025

\bibitem{memoryagentbench2026}
Hu, Y., et al.: MemoryAgentBench: Evaluating Memory in LLM Agents via
Incremental Multi-Turn Interactions. In: ICLR 2026

\bibitem{memoryarena2026}
MemoryArena: Evaluating Active Memory in Agentic Tasks. In: ICLR 2026

\bibitem{memoryr1_2025}
Memory-R1: RL Framework with Memory Manager for Agent Memory.
arXiv (2025)

\bibitem{agemem2026}
AgeMem: Agentic Memory with Unified LTM/STM Management via RL.
arXiv (2026)

\bibitem{gam2026}
Wu, Z., et al.: GAM: Hierarchical Graph Memory for LLM-based Agents.
In: ICLR 2026

\bibitem{memagents2026}
MemAgents: Memory for LLM-Based Agentic Systems. ICLR 2026 Workshop

\bibitem{minja2026}
Zeng, H., et al.: Memory Injection Attacks on LLM Agents via
Query-Only Interaction. In: MemAgents Workshop, ICLR 2026

\bibitem{wang2025mextra}
Wang, B., et al.: Unveiling Privacy Risks in LLM Agent Memory. In: ACL 2025

\bibitem{memarchitect2026}
Kumar, L.S., Ba, Y., Pan, R.: MemArchitect: A Policy Driven Memory
Governance Layer. arXiv:2603.18330 (2026)

\bibitem{ssgm2026}
Lam, C., et al.: Governing Evolving Memory in LLM Agents: Risks,
Mechanisms, and the SSGM Framework. arXiv:2603.11768 (2026)

\bibitem{mnemonic2026}
A Survey on the Security of Long-Term Memory in LLM Agents: Toward
Mnemonic Sovereignty. arXiv:2604.16548 (2026)

\bibitem{rezazadeh2025}
Rezazadeh, A., et al.: Collaborative Memory: Multi-User Memory Sharing
in LLM Agents with Dynamic Access Control. arXiv:2505.18279 (2025)

\bibitem{shi2025}
Shi, Z., et al.: Privacy-Enhancing Paradigms within Federated Multi-Agent
Systems. In: ICML 2025

\end{thebibliography}

\end{document}
